% =====================================================================
%  preambulo.tex
%  Paquetes, configuración de página y estilo para bloques de código.
%  Se carga desde main.tex con % =====================================================================
%  preambulo.tex
%  Paquetes, configuración de página y estilo para bloques de código.
%  Se carga desde main.tex con % =====================================================================
%  preambulo.tex
%  Paquetes, configuración de página y estilo para bloques de código.
%  Se carga desde main.tex con % =====================================================================
%  preambulo.tex
%  Paquetes, configuración de página y estilo para bloques de código.
%  Se carga desde main.tex con \input{preambulo}
% =====================================================================

% ---------------------------------------------------------------------
%  Codificación, idioma y fuentes
% ---------------------------------------------------------------------
\usepackage[T1]{fontenc}          % acentos y guiones correctos en español
\usepackage[utf8]{inputenc}
\usepackage{lmodern}              % versión vectorial de Computer Modern
\usepackage[spanish, provide=*]{babel}

% ---------------------------------------------------------------------
%  Geometría de página y encabezados
% ---------------------------------------------------------------------
\usepackage[a4paper, width=150mm, top=25mm, bottom=25mm, bindingoffset=6mm]{geometry}
\usepackage{fancyhdr}
\setlength{\headheight}{15.4pt}   % evita el warning de fancyhdr
\pagestyle{fancy}
\fancyhead{}
\fancyhead[LO,LE]{\leftmark}
\fancyfoot{}
\fancyfoot[CO,CE]{\thepage}
\renewcommand{\headrulewidth}{0.2pt}
\renewcommand{\footrulewidth}{0.2pt}

% ---------------------------------------------------------------------
%  Colores (antes que otros paquetes que puedan cargarlo)
% ---------------------------------------------------------------------
\usepackage[dvipsnames,table]{xcolor}

% ---------------------------------------------------------------------
%  Matemáticas y símbolos
% ---------------------------------------------------------------------
\usepackage{amsmath}
\usepackage{amssymb}
\usepackage{textcomp}
\usepackage{siunitx}

% ---------------------------------------------------------------------
%  Tablas
% ---------------------------------------------------------------------
\usepackage{array}
\usepackage{longtable}
\usepackage{multirow}
\usepackage{bigstrut}
\usepackage{booktabs}
\usepackage{tabularx}

% ---------------------------------------------------------------------
%  Figuras y páginas horizontales
% ---------------------------------------------------------------------
\usepackage{graphicx}
\usepackage{adjustbox}
\usepackage{float}
\usepackage{lscape}
\usepackage{pdflscape}

% ---------------------------------------------------------------------
%  Bibliografía (formato IEEE: citas numéricas entre corchetes)
% ---------------------------------------------------------------------
\usepackage{notoccite}
\usepackage[numbers, sort&compress]{natbib}
\usepackage[nottoc]{tocbibind}

% ---------------------------------------------------------------------
%  Código fuente: paquete listings + estilo "codigo"
% ---------------------------------------------------------------------
\usepackage{listings}

% Paleta inspirada en la imagen de referencia
\definecolor{codFondo}{HTML}{FAFAFB}      % fondo casi blanco
\definecolor{codMarco}{HTML}{D0D0D5}      % borde gris claro
\definecolor{codTexto}{HTML}{2B2B2B}      % variables e identificadores
\definecolor{codPalabra}{HTML}{5B8BC9}    % def, if, else, return...
\definecolor{codFuncion}{HTML}{4A7FC1}    % nombres de funciones (emph)
\definecolor{codNumero}{HTML}{6AA56A}     % números
\definecolor{codComentario}{HTML}{8C8C8C} % comentarios
\definecolor{codCadena}{HTML}{C1693C}     % cadenas de texto
\definecolor{codNumLinea}{HTML}{A0A0A0}   % números de línea

% Colorea los dígitos solo cuando están en código (no en comentarios ni cadenas)
\makeatletter
\newcommand{\codDigito}[1]{%
  \ifnum\lst@mode=\lst@Pmode\relax
    {\color{codNumero}#1}%
  \else
    #1%
  \fi
}
\makeatother

\lstdefinestyle{codigo}{
  language=Python,
  % --- Tipografía y colores ---
  basicstyle=\ttfamily\small\color{codTexto},
  identifierstyle=\color{codTexto},
  keywordstyle=\color{codPalabra},
  keywordstyle=[2]\color{codPalabra},
  emphstyle=\color{codFuncion},
  commentstyle=\color{codComentario},
  stringstyle=\color{codCadena},
  % --- Palabras reservadas adicionales ---
  morekeywords=[2]{True, False, None, self, nonlocal, async, await},
  % --- Marco y fondo ---
  backgroundcolor=\color{codFondo},
  frame=single,
  frameround=tttt,          % las cuatro esquinas redondeadas
  framerule=0.6pt,
  rulecolor=\color{codMarco},
  framesep=8pt,
  xleftmargin=9pt,
  xrightmargin=9pt,
  aboveskip=1em,
  belowskip=1em,
  % --- Espaciado y cortes de línea ---
  columns=fullflexible,
  keepspaces=true,
  showstringspaces=false,
  tabsize=4,
  breaklines=true,
  breakatwhitespace=true,
  postbreak=\mbox{\textcolor{codComentario}{$\hookrightarrow$}\space},
  upquote=true,
  % --- Números de línea (cambia a "left" si los quieres) ---
  numbers=none,
  numberstyle=\tiny\color{codNumLinea},
  % --- Acentos en español y coloreado de números ---
  extendedchars=true,
  literate=
    {á}{{\'a}}1 {é}{{\'e}}1 {í}{{\'i}}1 {ó}{{\'o}}1 {ú}{{\'u}}1
    {Á}{{\'A}}1 {É}{{\'E}}1 {Í}{{\'I}}1 {Ó}{{\'O}}1 {Ú}{{\'U}}1
    {ñ}{{\~n}}1 {Ñ}{{\~N}}1 {ü}{{\"u}}1 {Ü}{{\"U}}1
    {¿}{{?`}}1 {¡}{{!`}}1
    {0}{{\codDigito{0}}}1 {1}{{\codDigito{1}}}1 {2}{{\codDigito{2}}}1
    {3}{{\codDigito{3}}}1 {4}{{\codDigito{4}}}1 {5}{{\codDigito{5}}}1
    {6}{{\codDigito{6}}}1 {7}{{\codDigito{7}}}1 {8}{{\codDigito{8}}}1
    {9}{{\codDigito{9}}}1
}

% Estilo por defecto para todos los lstlisting
\lstset{style=codigo}

% Nombres en español para leyendas e índice de códigos
\renewcommand{\lstlistingname}{Código}
\renewcommand{\lstlistlistingname}{Índice de códigos}

% ---------------------------------------------------------------------
%  Enlaces (hyperref debe ir al final)
% ---------------------------------------------------------------------
\usepackage{url}
\usepackage[hidelinks]{hyperref}
% =====================================================================

% ---------------------------------------------------------------------
%  Codificación, idioma y fuentes
% ---------------------------------------------------------------------
\usepackage[T1]{fontenc}          % acentos y guiones correctos en español
\usepackage[utf8]{inputenc}
\usepackage{lmodern}              % versión vectorial de Computer Modern
\usepackage[spanish, provide=*]{babel}

% ---------------------------------------------------------------------
%  Geometría de página y encabezados
% ---------------------------------------------------------------------
\usepackage[a4paper, width=150mm, top=25mm, bottom=25mm, bindingoffset=6mm]{geometry}
\usepackage{fancyhdr}
\setlength{\headheight}{15.4pt}   % evita el warning de fancyhdr
\pagestyle{fancy}
\fancyhead{}
\fancyhead[LO,LE]{\leftmark}
\fancyfoot{}
\fancyfoot[CO,CE]{\thepage}
\renewcommand{\headrulewidth}{0.2pt}
\renewcommand{\footrulewidth}{0.2pt}

% ---------------------------------------------------------------------
%  Colores (antes que otros paquetes que puedan cargarlo)
% ---------------------------------------------------------------------
\usepackage[dvipsnames,table]{xcolor}

% ---------------------------------------------------------------------
%  Matemáticas y símbolos
% ---------------------------------------------------------------------
\usepackage{amsmath}
\usepackage{amssymb}
\usepackage{textcomp}
\usepackage{siunitx}

% ---------------------------------------------------------------------
%  Tablas
% ---------------------------------------------------------------------
\usepackage{array}
\usepackage{longtable}
\usepackage{multirow}
\usepackage{bigstrut}
\usepackage{booktabs}
\usepackage{tabularx}

% ---------------------------------------------------------------------
%  Figuras y páginas horizontales
% ---------------------------------------------------------------------
\usepackage{graphicx}
\usepackage{adjustbox}
\usepackage{float}
\usepackage{lscape}
\usepackage{pdflscape}

% ---------------------------------------------------------------------
%  Bibliografía (formato IEEE: citas numéricas entre corchetes)
% ---------------------------------------------------------------------
\usepackage{notoccite}
\usepackage[numbers, sort&compress]{natbib}
\usepackage[nottoc]{tocbibind}

% ---------------------------------------------------------------------
%  Código fuente: paquete listings + estilo "codigo"
% ---------------------------------------------------------------------
\usepackage{listings}

% Paleta inspirada en la imagen de referencia
\definecolor{codFondo}{HTML}{FAFAFB}      % fondo casi blanco
\definecolor{codMarco}{HTML}{D0D0D5}      % borde gris claro
\definecolor{codTexto}{HTML}{2B2B2B}      % variables e identificadores
\definecolor{codPalabra}{HTML}{5B8BC9}    % def, if, else, return...
\definecolor{codFuncion}{HTML}{4A7FC1}    % nombres de funciones (emph)
\definecolor{codNumero}{HTML}{6AA56A}     % números
\definecolor{codComentario}{HTML}{8C8C8C} % comentarios
\definecolor{codCadena}{HTML}{C1693C}     % cadenas de texto
\definecolor{codNumLinea}{HTML}{A0A0A0}   % números de línea

% Colorea los dígitos solo cuando están en código (no en comentarios ni cadenas)
\makeatletter
\newcommand{\codDigito}[1]{%
  \ifnum\lst@mode=\lst@Pmode\relax
    {\color{codNumero}#1}%
  \else
    #1%
  \fi
}
\makeatother

\lstdefinestyle{codigo}{
  language=Python,
  % --- Tipografía y colores ---
  basicstyle=\ttfamily\small\color{codTexto},
  identifierstyle=\color{codTexto},
  keywordstyle=\color{codPalabra},
  keywordstyle=[2]\color{codPalabra},
  emphstyle=\color{codFuncion},
  commentstyle=\color{codComentario},
  stringstyle=\color{codCadena},
  % --- Palabras reservadas adicionales ---
  morekeywords=[2]{True, False, None, self, nonlocal, async, await},
  % --- Marco y fondo ---
  backgroundcolor=\color{codFondo},
  frame=single,
  frameround=tttt,          % las cuatro esquinas redondeadas
  framerule=0.6pt,
  rulecolor=\color{codMarco},
  framesep=8pt,
  xleftmargin=9pt,
  xrightmargin=9pt,
  aboveskip=1em,
  belowskip=1em,
  % --- Espaciado y cortes de línea ---
  columns=fullflexible,
  keepspaces=true,
  showstringspaces=false,
  tabsize=4,
  breaklines=true,
  breakatwhitespace=true,
  postbreak=\mbox{\textcolor{codComentario}{$\hookrightarrow$}\space},
  upquote=true,
  % --- Números de línea (cambia a "left" si los quieres) ---
  numbers=none,
  numberstyle=\tiny\color{codNumLinea},
  % --- Acentos en español y coloreado de números ---
  extendedchars=true,
  literate=
    {á}{{\'a}}1 {é}{{\'e}}1 {í}{{\'i}}1 {ó}{{\'o}}1 {ú}{{\'u}}1
    {Á}{{\'A}}1 {É}{{\'E}}1 {Í}{{\'I}}1 {Ó}{{\'O}}1 {Ú}{{\'U}}1
    {ñ}{{\~n}}1 {Ñ}{{\~N}}1 {ü}{{\"u}}1 {Ü}{{\"U}}1
    {¿}{{?`}}1 {¡}{{!`}}1
    {0}{{\codDigito{0}}}1 {1}{{\codDigito{1}}}1 {2}{{\codDigito{2}}}1
    {3}{{\codDigito{3}}}1 {4}{{\codDigito{4}}}1 {5}{{\codDigito{5}}}1
    {6}{{\codDigito{6}}}1 {7}{{\codDigito{7}}}1 {8}{{\codDigito{8}}}1
    {9}{{\codDigito{9}}}1
}

% Estilo por defecto para todos los lstlisting
\lstset{style=codigo}

% Nombres en español para leyendas e índice de códigos
\renewcommand{\lstlistingname}{Código}
\renewcommand{\lstlistlistingname}{Índice de códigos}

% ---------------------------------------------------------------------
%  Enlaces (hyperref debe ir al final)
% ---------------------------------------------------------------------
\usepackage{url}
\usepackage[hidelinks]{hyperref}
% =====================================================================

% ---------------------------------------------------------------------
%  Codificación, idioma y fuentes
% ---------------------------------------------------------------------
\usepackage[T1]{fontenc}          % acentos y guiones correctos en español
\usepackage[utf8]{inputenc}
\usepackage{lmodern}              % versión vectorial de Computer Modern
\usepackage[spanish, provide=*]{babel}

% ---------------------------------------------------------------------
%  Geometría de página y encabezados
% ---------------------------------------------------------------------
\usepackage[a4paper, width=150mm, top=25mm, bottom=25mm, bindingoffset=6mm]{geometry}
\usepackage{fancyhdr}
\setlength{\headheight}{15.4pt}   % evita el warning de fancyhdr
\pagestyle{fancy}
\fancyhead{}
\fancyhead[LO,LE]{\leftmark}
\fancyfoot{}
\fancyfoot[CO,CE]{\thepage}
\renewcommand{\headrulewidth}{0.2pt}
\renewcommand{\footrulewidth}{0.2pt}

% ---------------------------------------------------------------------
%  Colores (antes que otros paquetes que puedan cargarlo)
% ---------------------------------------------------------------------
\usepackage[dvipsnames,table]{xcolor}

% ---------------------------------------------------------------------
%  Matemáticas y símbolos
% ---------------------------------------------------------------------
\usepackage{amsmath}
\usepackage{amssymb}
\usepackage{textcomp}
\usepackage{siunitx}

% ---------------------------------------------------------------------
%  Tablas
% ---------------------------------------------------------------------
\usepackage{array}
\usepackage{longtable}
\usepackage{multirow}
\usepackage{bigstrut}
\usepackage{booktabs}
\usepackage{tabularx}

% ---------------------------------------------------------------------
%  Figuras y páginas horizontales
% ---------------------------------------------------------------------
\usepackage{graphicx}
\usepackage{adjustbox}
\usepackage{float}
\usepackage{lscape}
\usepackage{pdflscape}

% ---------------------------------------------------------------------
%  Bibliografía (formato IEEE: citas numéricas entre corchetes)
% ---------------------------------------------------------------------
\usepackage{notoccite}
\usepackage[numbers, sort&compress]{natbib}
\usepackage[nottoc]{tocbibind}

% ---------------------------------------------------------------------
%  Código fuente: paquete listings + estilo "codigo"
% ---------------------------------------------------------------------
\usepackage{listings}

% Paleta inspirada en la imagen de referencia
\definecolor{codFondo}{HTML}{FAFAFB}      % fondo casi blanco
\definecolor{codMarco}{HTML}{D0D0D5}      % borde gris claro
\definecolor{codTexto}{HTML}{2B2B2B}      % variables e identificadores
\definecolor{codPalabra}{HTML}{5B8BC9}    % def, if, else, return...
\definecolor{codFuncion}{HTML}{4A7FC1}    % nombres de funciones (emph)
\definecolor{codNumero}{HTML}{6AA56A}     % números
\definecolor{codComentario}{HTML}{8C8C8C} % comentarios
\definecolor{codCadena}{HTML}{C1693C}     % cadenas de texto
\definecolor{codNumLinea}{HTML}{A0A0A0}   % números de línea

% Colorea los dígitos solo cuando están en código (no en comentarios ni cadenas)
\makeatletter
\newcommand{\codDigito}[1]{%
  \ifnum\lst@mode=\lst@Pmode\relax
    {\color{codNumero}#1}%
  \else
    #1%
  \fi
}
\makeatother

\lstdefinestyle{codigo}{
  language=Python,
  % --- Tipografía y colores ---
  basicstyle=\ttfamily\small\color{codTexto},
  identifierstyle=\color{codTexto},
  keywordstyle=\color{codPalabra},
  keywordstyle=[2]\color{codPalabra},
  emphstyle=\color{codFuncion},
  commentstyle=\color{codComentario},
  stringstyle=\color{codCadena},
  % --- Palabras reservadas adicionales ---
  morekeywords=[2]{True, False, None, self, nonlocal, async, await},
  % --- Marco y fondo ---
  backgroundcolor=\color{codFondo},
  frame=single,
  frameround=tttt,          % las cuatro esquinas redondeadas
  framerule=0.6pt,
  rulecolor=\color{codMarco},
  framesep=8pt,
  xleftmargin=9pt,
  xrightmargin=9pt,
  aboveskip=1em,
  belowskip=1em,
  % --- Espaciado y cortes de línea ---
  columns=fullflexible,
  keepspaces=true,
  showstringspaces=false,
  tabsize=4,
  breaklines=true,
  breakatwhitespace=true,
  postbreak=\mbox{\textcolor{codComentario}{$\hookrightarrow$}\space},
  upquote=true,
  % --- Números de línea (cambia a "left" si los quieres) ---
  numbers=none,
  numberstyle=\tiny\color{codNumLinea},
  % --- Acentos en español y coloreado de números ---
  extendedchars=true,
  literate=
    {á}{{\'a}}1 {é}{{\'e}}1 {í}{{\'i}}1 {ó}{{\'o}}1 {ú}{{\'u}}1
    {Á}{{\'A}}1 {É}{{\'E}}1 {Í}{{\'I}}1 {Ó}{{\'O}}1 {Ú}{{\'U}}1
    {ñ}{{\~n}}1 {Ñ}{{\~N}}1 {ü}{{\"u}}1 {Ü}{{\"U}}1
    {¿}{{?`}}1 {¡}{{!`}}1
    {0}{{\codDigito{0}}}1 {1}{{\codDigito{1}}}1 {2}{{\codDigito{2}}}1
    {3}{{\codDigito{3}}}1 {4}{{\codDigito{4}}}1 {5}{{\codDigito{5}}}1
    {6}{{\codDigito{6}}}1 {7}{{\codDigito{7}}}1 {8}{{\codDigito{8}}}1
    {9}{{\codDigito{9}}}1
}

% Estilo por defecto para todos los lstlisting
\lstset{style=codigo}

% Nombres en español para leyendas e índice de códigos
\renewcommand{\lstlistingname}{Código}
\renewcommand{\lstlistlistingname}{Índice de códigos}

% ---------------------------------------------------------------------
%  Enlaces (hyperref debe ir al final)
% ---------------------------------------------------------------------
\usepackage{url}
\usepackage[hidelinks]{hyperref}
% =====================================================================

% ---------------------------------------------------------------------
%  Codificación, idioma y fuentes
% ---------------------------------------------------------------------
\usepackage[T1]{fontenc}          % acentos y guiones correctos en español
\usepackage[utf8]{inputenc}
\usepackage{lmodern}              % versión vectorial de Computer Modern
\usepackage[spanish, provide=*]{babel}

% ---------------------------------------------------------------------
%  Geometría de página y encabezados
% ---------------------------------------------------------------------
\usepackage[a4paper, width=150mm, top=25mm, bottom=25mm, bindingoffset=6mm]{geometry}
\usepackage{fancyhdr}
\setlength{\headheight}{15.4pt}   % evita el warning de fancyhdr
\pagestyle{fancy}
\fancyhead{}
\fancyhead[LO,LE]{\leftmark}
\fancyfoot{}
\fancyfoot[CO,CE]{\thepage}
\renewcommand{\headrulewidth}{0.2pt}
\renewcommand{\footrulewidth}{0.2pt}

% ---------------------------------------------------------------------
%  Colores (antes que otros paquetes que puedan cargarlo)
% ---------------------------------------------------------------------
\usepackage[dvipsnames,table]{xcolor}

% ---------------------------------------------------------------------
%  Matemáticas y símbolos
% ---------------------------------------------------------------------
\usepackage{amsmath}
\usepackage{amssymb}
\usepackage{textcomp}
\usepackage{siunitx}

% ---------------------------------------------------------------------
%  Tablas
% ---------------------------------------------------------------------
\usepackage{array}
\usepackage{longtable}
\usepackage{multirow}
\usepackage{bigstrut}
\usepackage{booktabs}
\usepackage{tabularx}

% ---------------------------------------------------------------------
%  Figuras y páginas horizontales
% ---------------------------------------------------------------------
\usepackage{graphicx}
\usepackage{adjustbox}
\usepackage{float}
\usepackage{lscape}
\usepackage{pdflscape}

% ---------------------------------------------------------------------
%  Bibliografía (formato IEEE: citas numéricas entre corchetes)
% ---------------------------------------------------------------------
\usepackage{notoccite}
\usepackage[numbers, sort&compress]{natbib}
\usepackage[nottoc]{tocbibind}

% ---------------------------------------------------------------------
%  Código fuente: paquete listings + estilo "codigo"
% ---------------------------------------------------------------------
\usepackage{listings}

% Paleta inspirada en la imagen de referencia
\definecolor{codFondo}{HTML}{FAFAFB}      % fondo casi blanco
\definecolor{codMarco}{HTML}{D0D0D5}      % borde gris claro
\definecolor{codTexto}{HTML}{2B2B2B}      % variables e identificadores
\definecolor{codPalabra}{HTML}{5B8BC9}    % def, if, else, return...
\definecolor{codFuncion}{HTML}{4A7FC1}    % nombres de funciones (emph)
\definecolor{codNumero}{HTML}{6AA56A}     % números
\definecolor{codComentario}{HTML}{8C8C8C} % comentarios
\definecolor{codCadena}{HTML}{C1693C}     % cadenas de texto
\definecolor{codNumLinea}{HTML}{A0A0A0}   % números de línea

% Colorea los dígitos solo cuando están en código (no en comentarios ni cadenas)
\makeatletter
\newcommand{\codDigito}[1]{%
  \ifnum\lst@mode=\lst@Pmode\relax
    {\color{codNumero}#1}%
  \else
    #1%
  \fi
}
\makeatother

\lstdefinestyle{codigo}{
  language=Python,
  % --- Tipografía y colores ---
  basicstyle=\ttfamily\small\color{codTexto},
  identifierstyle=\color{codTexto},
  keywordstyle=\color{codPalabra},
  keywordstyle=[2]\color{codPalabra},
  emphstyle=\color{codFuncion},
  commentstyle=\color{codComentario},
  stringstyle=\color{codCadena},
  % --- Palabras reservadas adicionales ---
  morekeywords=[2]{True, False, None, self, nonlocal, async, await},
  % --- Marco y fondo ---
  backgroundcolor=\color{codFondo},
  frame=single,
  frameround=tttt,          % las cuatro esquinas redondeadas
  framerule=0.6pt,
  rulecolor=\color{codMarco},
  framesep=8pt,
  xleftmargin=9pt,
  xrightmargin=9pt,
  aboveskip=1em,
  belowskip=1em,
  % --- Espaciado y cortes de línea ---
  columns=fullflexible,
  keepspaces=true,
  showstringspaces=false,
  tabsize=4,
  breaklines=true,
  breakatwhitespace=true,
  postbreak=\mbox{\textcolor{codComentario}{$\hookrightarrow$}\space},
  upquote=true,
  % --- Números de línea (cambia a "left" si los quieres) ---
  numbers=none,
  numberstyle=\tiny\color{codNumLinea},
  % --- Acentos en español y coloreado de números ---
  extendedchars=true,
  literate=
    {á}{{\'a}}1 {é}{{\'e}}1 {í}{{\'i}}1 {ó}{{\'o}}1 {ú}{{\'u}}1
    {Á}{{\'A}}1 {É}{{\'E}}1 {Í}{{\'I}}1 {Ó}{{\'O}}1 {Ú}{{\'U}}1
    {ñ}{{\~n}}1 {Ñ}{{\~N}}1 {ü}{{\"u}}1 {Ü}{{\"U}}1
    {¿}{{?`}}1 {¡}{{!`}}1
    {0}{{\codDigito{0}}}1 {1}{{\codDigito{1}}}1 {2}{{\codDigito{2}}}1
    {3}{{\codDigito{3}}}1 {4}{{\codDigito{4}}}1 {5}{{\codDigito{5}}}1
    {6}{{\codDigito{6}}}1 {7}{{\codDigito{7}}}1 {8}{{\codDigito{8}}}1
    {9}{{\codDigito{9}}}1
}

% Estilo por defecto para todos los lstlisting
\lstset{style=codigo}

% Nombres en español para leyendas e índice de códigos
\renewcommand{\lstlistingname}{Código}
\renewcommand{\lstlistlistingname}{Índice de códigos}

% ---------------------------------------------------------------------
%  Enlaces (hyperref debe ir al final)
% ---------------------------------------------------------------------
\usepackage{url}
\usepackage[hidelinks]{hyperref}